\chapter{Geometry of the Expanding Universe}

The point is the expanding universe and its geometry. The aim is to describe what it means for space to expand before asking what drives the expansion. Newtonian \(F=ma\) cosmology comes later, and only after that do geometry and dynamics get assembled into the full Friedmann--Robertson--Walker framework.

\section{Comoving Coordinates on an Expanding Line}

Start with the simplest possible example: one-dimensional space, not spacetime. Imagine an infinite line with equally spaced galaxies on it. The coordinate \(x\) does not yet measure a physical distance. It labels the galaxies. One may think of \(x\) as the name of a galaxy, and that name stays attached to the galaxy even if the whole arrangement stretches or contracts.

Let the physical distance from the galaxy called \(x=0\) to the galaxy called \(x=1\) be \(a\). Then the physical distance between any two galaxies is
\begin{equation}
d=a\,\Delta x.
\end{equation}
The coordinate difference \(\Delta x\) is only the difference between labels. The factor \(a\) converts that label difference into an actual distance measured by meter sticks.

Now imagine the line as an elastic thread being stretched. The labels stay with the galaxies, but the distance between neighboring galaxies changes with time. The same relation becomes
\begin{equation}
d(t)=a(t)\,\Delta x.
\end{equation}
At this stage homogeneity means simply that neighboring galaxies are equally spaced everywhere, even though that common spacing may vary with time.

\section{Hubble's Law and the Absence of a Center}

To find the relative velocity of two fixed galaxies, keep the pair fixed so that \(\Delta x\) does not change with time. Differentiating gives
\begin{align}
\dot d &= \dot a\,\Delta x \\
       &= \frac{\dot a}{a}\,a\,\Delta x \\
       &= \frac{\dot a}{a}\,d.
\end{align}
If we define
\begin{equation}
H(t)\equiv \frac{\dot a}{a},
\end{equation}
then
\begin{equation}
\dot d = H(t)\,d.
\end{equation}
This is Hubble's law.

The historical phrase ``Hubble constant'' is misleading. The quantity is constant in space, not in time. At a given moment it is the same for every galaxy, but in general it changes as the universe evolves. It is better thought of as the Hubble parameter.

The law says that more distant galaxies recede faster. It does not single out a center. Nothing in the argument privileges \(x=0\), or any other point. Every observer sees the same rule. The expanding universe in this model is therefore not a blast from one distinguished location.

Redshift enters immediately. If two galaxies exchange light signals while moving apart, the received light is reddened. The greater the recession speed, the greater the redshift. Spectral lines therefore tell us about recession speeds. Distances are harder to measure, and Hubble's early estimates were rough. His first plot was not a perfect straight line but a scattered set of points with large uncertainties. Even so, the geometric principle survived the noise: recession speed is proportional to distance.

\begin{classroomqa}
\audiencequestion{If the galaxies are far enough apart, does Hubble's law make them move faster than the speed of light?}\lecturetimestamp{00:20:23}
\lecturerresponse{Yes, in this recession sense it does. What you are saying is correct. We will come back to that for sure. It is one of the most important things.}
\end{classroomqa}

\section{Higher Dimensions, Isotropy, and the Time Dependence of \(a(t)\)}

Nothing essential changes when more spatial dimensions are added. In two dimensions imagine a square grid of galaxies labeled by \((x,y)\), and in three dimensions by \((x,y,z)\). The coordinates are again comoving labels. They move with the galaxies instead of measuring physical distance directly.

Why should the spacing in the \(x\)-direction be the same as the spacing in the \(y\)-direction? The answer is observational. If galaxies were systematically more crowded in one direction than another, the sky would look different in different directions. On sufficiently large scales it does not. That directional sameness is isotropy.

In two dimensions the distance between two galaxies is
\begin{equation}
d=a\sqrt{(\Delta x)^2+(\Delta y)^2},
\end{equation}
and in three dimensions
\begin{equation}
d=a\sqrt{(\Delta x)^2+(\Delta y)^2+(\Delta z)^2}.
\end{equation}
For a fixed pair of galaxies the coordinate differences remain constant, so differentiating again gives
\begin{equation}
\dot d=\frac{\dot a}{a}\,d.
\end{equation}
Hubble's law is unchanged.

Homogeneity means sameness from place to place. Isotropy means sameness from direction to direction. Neither is exact on small scales. Looking through the Milky Way is different from looking out of its plane, and individual galaxies are not uniformly spread at short distances. But after averaging over many galaxies, the large-scale universe looks approximately the same everywhere and in every direction.

\begin{classroomqa}
\audiencequestion{If velocity is proportional to distance, and the distance is increasing, does that mean a fixed pair of galaxies must be accelerating away from each other?}\lecturetimestamp{00:40:22}
\lecturerresponse{No. Hubble's law says that at a given moment the recession velocity is proportional to the distance. It does not by itself say whether the velocity of one fixed pair increases or decreases with time. That depends on the detailed time dependence of \(a(t)\). If \(a(t)\propto t\), neighboring galaxies separate at constant speed. If \(a(t)\propto t^2\), the recession speed of a fixed pair increases with time. If \(a(t)\propto \log t\), it decreases with time. The law \(\dot d=(\dot a/a)d\) is always true as a function of distance, but acceleration and deceleration are separate questions.}
\end{classroomqa}

For a long time cosmologists expected the universe to be decelerating, since gravity should slow an initial expansion the way gravity slows an object thrown upward. Here the point is only that Hubble's law by itself does not determine the behavior of \(\dot a\) or \(\ddot a\).

\begin{classroomqa}
\audiencequestion{How do you set the initial value of the scale factor \(a\)?}\lecturetimestamp{00:46:18}
\lecturerresponse{Its overall normalization is arbitrary. What matters physically is not the numerical size of \(a\) by itself, but how it changes. If one defines another scale factor \(b=3a\), then \(\dot b/b=\dot a/a\). Rescaling \(a\) by a constant changes conventions or units, not the Hubble parameter.}
\end{classroomqa}

\section{From the Spatial Metric to the Spacetime Metric}

Now replace \(\Delta x\) by \(dx\). The physical distance between neighboring points is
\begin{equation}
ds=a\,dx,
\end{equation}
so
\begin{equation}
ds^2=a^2\,dx^2.
\end{equation}
This is the spatial metric in one dimension. The coefficient of \(dx^2\) is
\begin{equation}
g_{xx}=a^2.
\end{equation}
An expanding universe can therefore be described by saying that the spatial part of the metric depends on time.

To pass from space to spacetime, use the relativistic interval. In one spatial dimension,
\begin{equation}
d\tau^2=dt^2-a^2(t)\,dx^2.
\end{equation}
For a comoving observer, \(dx=0\), so \(d\tau=dt\). Ordinary cosmic time is the proper time measured by clocks moving with the galaxies.

The blackboard picture at this point is useful because it puts the comoving worldlines, the line element, and the metric notation together.
\begin{figure}[H]
\centering
\includegraphics[width=0.82\textwidth]{lecture_01_figure_02.png}
\caption{Comoving worldlines in an expanding one-dimensional universe together with the corresponding spacetime interval.}
\end{figure}

In three spatial dimensions the same pattern gives the cautious extension\footnote{The board is partly hidden around 00:56:38--00:58:51. This three-dimensional form follows the spoken extension from one spatial dimension to three together with the visible diagonal pattern on the board.}
\begin{equation}
d\tau^2=dt^2-a^2(t)\bigl(dx^2+dy^2+dz^2\bigr),
\end{equation}
with diagonal metric
\begin{equation}
g_{\mu\nu}=\mathrm{diag}\bigl(1,-a^2,-a^2,-a^2\bigr).
\end{equation}
The off-diagonal terms vanish,
\begin{equation}
g_{tx}=g_{ty}=g_{tz}=0.
\end{equation}

Homogeneity means that the metric coefficients may depend on time but not on spatial position. Isotropy means that the three spatial directions enter in the same way.

\begin{classroomqa}
\audiencequestion{Are we not supposed to think of \(dx\) as a distance?}\lecturetimestamp{01:01:18}
\lecturerresponse{No. \(dx\) is not the distance. \(a\,dx\) is the distance. \(dx\) is just a label interval. A useful way to think about it is to imagine that, at some moment, someone draws a coordinate mesh through the galaxies and then requires the mesh to follow them. The coordinate separation stays fixed by construction, while the physical distance changes. One could put the units of distance into \(x\) instead of into \(a\); that is a convention. What matters is the product \(a\,dx\).}
\end{classroomqa}

\begin{classroomqa}
\audiencequestion{The metric coefficient of \(dt^2\) is equal to \(1\). Is that just by convention also?}\lecturetimestamp{01:03:24}
\lecturerresponse{Yes. There is no physics in that by itself. If the coefficient depended only on time, one could change the time variable so that it became just \(dt^2\). Likewise one may choose units with \(c=1\). What is not a convention here is that the coefficient does not depend on position, although it may depend on time.}
\end{classroomqa}

\begin{classroomqa}
\audiencequestion{If you put this into Einstein's field equations, what do you get for the energy-momentum tensor?}\lecturetimestamp{01:05:59}
\lecturerresponse{That is the next kind of question, but it belongs to dynamics and not to the present discussion. One can put this metric into Einstein's equations and ask what energy-momentum distribution would make it a solution. That is very important, but for the moment the subject is geometry, not dynamics.}
\end{classroomqa}

\section{What the Horizon Hides}

The large-scale universe is treated as homogeneous and isotropic only as an approximation. It is certainly not homogeneous on the scale of stars, planets, or even individual galaxies. It is not exactly isotropic in every local observation either.

A farmer's field gives the right scale dependence. On one scale it looks regular; on a much smaller scale it plainly does not; on a scale larger than the whole field it again ceases to look homogeneous. Cosmology works the same way. After averaging over many galaxies, the large-scale distribution looks roughly uniform. On sufficiently small scales it does not.

The universe is bigger than the part we can see, and bigger than the part we can ever see. On the scales that can actually be probed, it is isotropic and homogeneous apart from local variations. Beyond that there is no direct observational guarantee.

\begin{classroomqa}
\audiencequestion{You said we will never really directly find out. Does that imply that there might be an indirect way to find out?}\lecturetimestamp{01:10:34}
\lecturerresponse{Perhaps. Most thinking these days is that on big enough scales the universe is not homogeneous, and my own bet is the same. But it is a bet I will never win and never lose, because for basic fundamental reasons we will not be able to see much farther than we can see now. That is the existence of a horizon.}
\end{classroomqa}

\begin{classroomqa}
\audiencequestion{Couldn't you use gravity to look for masses beyond what you can see?}\lecturetimestamp{01:11:53}
\lecturerresponse{No. A distant mass can simply put us into free fall. Sitting in a closed room, we are in the gravitational field of the Sun and orbiting it, but without seeing outside there is no easy way to tell that the Sun is there. Only tidal forces reveal it, and if the source is far enough away those tidal effects can be too small to detect.}
\end{classroomqa}

The horizon is therefore an ultimate barrier to direct knowledge. He even says, with virtual certainty, that the universe is at least a thousand times larger in volume than the region we can ever see, and leaves the reason for that claim for later.

\section{Bounded Without Boundary}

Now come to closed geometries. The infinite line is open and unbounded. A line segment is bounded in the sense of having only a finite extent. Bounded does not mean a boundary, but finite extent.

A cube with walls will not do. It violates homogeneity, because somebody near a wall is not equivalent to somebody far from a wall, and it violates isotropy, because even somebody at the center sees different things in different directions.

The interior of a circle or ball also fails. Somebody at the center sees the same thing in every direction, but that symmetry belongs only to one special point. Move away from the center and the wall is near in one direction and far in another. The space is then neither homogeneous nor isotropic for generic observers.

On very large scales the working assumption is that the universe is homogeneous and isotropic. This is the cosmological principle. It is useful, but it is not protected from revision.

\begin{classroomqa}
\audiencequestion{Does isotropy everywhere imply homogeneity?}\lecturetimestamp{01:18:57}
\lecturerresponse{Yes. If every observer sees the same thing in every direction, that implies homogeneity. But the converse is false. Homogeneity by itself does not imply isotropy. A universe can be homogeneous and still have a preferred direction, for example if it were filled with arrows all pointing the same way.}
\end{classroomqa}

The right bounded model is the surface of a featureless sphere. A two-dimensional creature confined to that surface finds every place locally like every other place and every direction like every other direction, even though the world is finite. That is bounded without boundary.

\section{A Closed One-Dimensional Universe on a Circle}

To keep the algebra simple, reduce the sphere example to one spatial dimension: life on a circle. In one dimension isotropy means that left and right are equivalent. A circle has exactly that feature. Choose one point and call it \(\theta=0\). Every other point is labeled by an angular coordinate \(\theta\), measured in radians. A point just to one side of the origin can be described either by a small negative angle or by an angle just short of \(2\pi\). That periodicity is part of the geometry.

The inhabitants of the circle have access to the circumference, not to an external radius or center. They can measure the total distance around by marching along the circle until they return to the same point. If one draws the circle in an ambient plane, one may also introduce a radius-like quantity, but that belongs to the embedding picture, not to the intrinsic geometry they measure.

The blackboard picture is useful here because it keeps the angular labels and the total circumference on the same drawing.
\begin{figure}[H]
\centering
\includegraphics[width=0.82\textwidth]{lecture_01_figure_03.png}
\caption{Angular description of a closed one-dimensional universe, with total circumference \(2\pi a\).}
\end{figure}

A convenient convention is\footnote{Around 01:26:40--01:29:20 the notation shifts briefly between taking the whole circumference itself as the scale factor and writing it as \(2\pi a\). The later convention is used here because it is the one carried into the arc-length formula.}
\begin{equation}
L=2\pi a,
\end{equation}
where \(L\) is the total circumference.

For a small angular separation \(\Delta\theta\), the physical distance is
\begin{equation}
d=a\,\Delta\theta.
\end{equation}
Passing to infinitesimals gives
\begin{equation}
ds=a\,d\theta,
\end{equation}
and therefore
\begin{equation}
ds^2=a^2\,d\theta^2.
\end{equation}
This is the local metric on the circle.

If \(a\) depends on time, the closed one-dimensional universe expands or contracts. For a fixed pair of galaxies,
\begin{equation}
d(t)=a(t)\,\Delta\theta,
\end{equation}
so
\begin{equation}
\dot d=\dot a\,\Delta\theta
      =\frac{\dot a}{a}\,d.
\end{equation}
Hubble's law survives here as well, at least for pairs that are not so far apart that the other way around the circle becomes ambiguous.

The spacetime metric becomes
\begin{equation}
d\tau^2=dt^2-a^2(t)\,d\theta^2.
\end{equation}
If \(a\) is constant, the spacetime picture is a cylinder. If \(a\propto t\), it is a cone. If \(a(t)\) changes more irregularly, the picture becomes vase-like.

Run the picture backward and the galaxies get closer and closer together. If there were \(360\) galaxies around the circle, they would be well spread out at late times, very close together earlier, then closer still. Eventually they would not be separate galaxies anymore, and earlier still not even separate atoms. But there is still no special point in space.

That is the cleanest answer to the question of where the Big Bang happened in the model. As \(a(t)\to 0\), every separation \(a(t)\,\Delta\theta\) goes to zero. The whole universe shrinks everywhere at once. The galaxies are not flying away from one privileged location in an external space. The geometry itself is changing scale. In that sense the answer is: everywhere.

\begin{classroomqa}
\audiencequestion{You said the spacetime is curved. Is it curved if that were a cylinder?}\lecturetimestamp{01:38:04}
\lecturerresponse{No. If it were a cylinder, it would not be curved. That is the case where \(a\) is constant in time. It looks curved only because it is drawn in a higher-dimensional space. If one cuts it open, it can be laid flat without stretching.}
\end{classroomqa}

From there the discussion turns to the cone. If \(a\propto t\), the spacetime picture is a cone. Away from its tip the cone is also flat in the same intrinsic sense: it can be opened up into a plane. The tip is different. It is a conical singularity.

The next step is the geometry of the spherical universe and then the dynamics that determine how \(a(t)\) grows or shrinks.
